% SPDX-License-Identifier: CC-BY-4.0
\documentclass[11pt,letterpaper]{article}

\usepackage[T1]{fontenc}
\usepackage{amsmath,amssymb}
\usepackage{tikz}
\usepackage{pgfplots}
\pgfplotsset{compat=1.18}
\usepackage{algorithm}
\usepackage{algpseudocode}
\usepackage{booktabs}
\usepackage{enumitem}
\usepackage[hidelinks]{hyperref}

\newcommand{\N}{\mathbb{N}}
\newcommand{\Z}{\mathbb{Z}}
\newcommand{\Q}{\mathbb{Q}}
\newcommand{\R}{\mathbb{R}}
\newcommand{\C}{\mathbb{C}}
\newcommand{\Map}[3]{#1 \colon #2 \to #3}
\newcommand{\Mod}[1]{#1\text{-}\mathrm{Mod}}
\newcommand{\GroupHom}{\mathrm{Hom}_{\mathbf{Grp}}}
\newcommand{\RingHom}{\mathrm{Hom}_{\mathbf{Ring}}}

\newtheorem{theorem}{Theorem}
\newtheorem{lemma}[theorem]{Lemma}
\newtheorem{corollary}[theorem]{Corollary}
\newtheorem{definition}[theorem]{Definition}
\newtheorem{remark}[theorem]{Remark}
\newenvironment{proof}{\noindent\textit{Proof.}}{\hfill\(\square\)}

\title{Explicit cycles and the power-of-\(d\) intersection}
\author{}
\date{\today}

\begin{document}
\maketitle

\section*{Explicit cycles and the power-of-\(d\) intersection}

Throughout, fix natural numbers \(d\ge2\), \(q\ge1\) and \(c\ge0\) and define
the raw map
\[
 T_{q,c,d}(n)=\begin{cases} n/d, & d\mid n,\\ qn+c, & d\nmid n,\end{cases}
 \qquad n\in\N .
\]
Division has precedence, \(T^{k}\) denotes \(k\) instructions of \(T\), so
\(T^{0}\) is the identity and \(T^{k+1}=T\circ T^{k}\). All states are natural
numbers, and a state is \emph{positive} when it is nonzero. Every division
below is exact where it is written.

The document proves three families of cycles, one dichotomy at the divisor
\(2\), and one intersection statement. The families are the geometric-offset
slice \(c=3(8^{h}-1)/7\), \(q=2c+1\), \(d=2\); the full-residue-access slice
\(c=(d-1)(d^{k-1}-1)\), \(q=(d-1)d^{k-1}+1\); and the divisor-two slices with
\(q=2c+1\), which also carry the dichotomy; the fourth and last statement is
the intersection slice \(q=d+1\), \(c=d-1\) for \(d\ge3\).

\section{The raw map and its iterates}

\begin{lemma}[Branch rules]\label{Lem:Branch}
For every \(n\),
\[
 T_{q,c,d}(n)=n/d \quad\text{if } d\mid n,
 \qquad
 T_{q,c,d}(n)=qn+c \quad\text{if } d\nmid n .
\]
In particular \(T_{q,c,d}(0)=0\) for all \(d\), \(q\), \(c\).
\end{lemma}

\begin{proof}
Both identities are the definition of \(T_{q,c,d}\), split by the decidable
predicate \(d\mid n\). For the last claim, \(d\mid0\), so the division branch
applies and \(0/d=0\).
\end{proof}

\begin{lemma}[Iterate decomposition]\label{Lem:IterateSplit}
Let \(f\) be a self-map of a type, \(k\in\N\) and \(x\) a state. Then
\[
 f^{k+1}(x)=f\bigl(f^{k}(x)\bigr)=f^{k}\bigl(f(x)\bigr),
 \qquad\text{and}\qquad
 f^{m+n}(x)=f^{n}\bigl(f^{m}(x)\bigr) .
\]
\end{lemma}

\begin{proof}
The first identity is the recurrence \(f^{k+1}=f\circ f^{k}\), the second its
mirror \(f^{k+1}=f^{k}\circ f\), which holds because powers of one map
commute; both follow by induction on \(k\). The third is the addition rule for
iterates, by induction on \(n\).
\end{proof}

\begin{lemma}[Divisions]\label{Lem:Divisions}
Let \(d\ge2\). Then for all \(m\) and \(j\),
\[
 T_{q,c,d}^{\,j}\bigl(d^{j}m\bigr)=m ,
 \qquad\text{and in particular}\qquad
 T_{q,c,d}^{\,j}\bigl(d^{j}\bigr)=1 .
\]
\end{lemma}

\begin{proof}
Induction on \(j\). For \(j=0\) both sides are \(m\), respectively \(1\). For
the step, \(d\) divides \(d^{j+1}m\), so the first instruction divides:
\(T(d^{j+1}m)=d^{j+1}m/d=d^{j}m\), and the induction hypothesis applies to the
remaining \(j\) instructions.
\end{proof}

\begin{lemma}[Bounded orbits repeat]\label{Lem:Pigeonhole}
If \(\;T_{q,c,d}^{\,k}(n)\le M\) for every \(k\), then there are \(k\) and
\(p>0\) with
\[
 T_{q,c,d}^{\,k+p}(n)=T_{q,c,d}^{\,k}(n) .
\]
\end{lemma}

\begin{proof}
The \(M+2\) states \(T^{0}(n),\dots,T^{M+1}(n)\) all lie in the \(M+1\) values
\(\{0,\dots,M\}\), so two of them agree at distinct indices \(i<j\). Set
\(k=i\) and \(p=j-i>0\). Then
\(T^{k+p}(n)=T^{\,j-i}(T^{\,i}(n))=T^{\,j}(n)=T^{\,i}(n)=T^{k}(n)\)
by the addition rule of Lemma~\ref{Lem:IterateSplit}.
\end{proof}

\section{The divisor-two slices}

\begin{lemma}[Odd states take the affine branch]\label{Lem:TwoOdd}
Let \(n\) be odd. Then \(T_{q,c,2}(n)=qn+c\), and with \(q=2c+1\),
\[
 T_{2c+1,c,2}(n)=(2c+1)n+c .
\]
If \(c\) is even then this value is odd again, and if in addition \(1\le c\)
then it exceeds \(n\).
\end{lemma}

\begin{proof}
An odd state is not divisible by \(2\), so the affine branch applies by
Lemma~\ref{Lem:Branch}, which gives the first two assertions. For the parity,
write \(c=2e\); then \((2c+1)n+c=2n(c+e)+n\), a sum of an even number and an
odd number, hence odd. For the growth,
\((2c+1)n+c=n+c(2n+1)\) with \(c(2n+1)>0\) because \(1\le c\) and \(0<n\).
\end{proof}

\begin{theorem}[Even offsets escape]\label{Thm:EvenEscape}
Let \(c\) be even with \(1\le c\), and let \(n\ge1\). Then every positive seed
of the slice \(T_{2c+1,c,2}\) reaches an odd state, and from there the orbit is
odd at every later time and passes every bound: for every \(M\) there is \(k\)
with \(M<T_{2c+1,c,2}^{\,k}(n)\).
\end{theorem}

\begin{proof}
For the reach, use strong induction on \(n\). If \(n\) is odd, the time \(0\)
already gives an odd state. If \(n\) is even, then \(2\mid n\), so
\(T(n)=n/2\), which is positive because \(1\le n\), and \(n/2<n\); the
induction hypothesis for \(n/2\) gives a time \(j\) with \(T^{\,j}(n/2)\) odd,
hence \(T^{\,j+1}(n)\) is odd.

For the escape, fix \(j\) with \(y=T^{\,j}(n)\) odd. By
Lemma~\ref{Lem:TwoOdd} the parity of odd states is preserved, so
\(T^{\,t}(y)\) is odd for every \(t\), and each step increases the state by at
least one. Hence \(T^{\,t}(y)\ge y+t\) by induction on \(t\), and choosing
\(t=M\) gives \(T^{\,j+M}(n)=T^{\,M}(y)\ge y+M>M\).
\end{proof}

\begin{lemma}[Odd offsets preserve the gcd]\label{Lem:OddGcd}
Let \(c\) be odd. Then for every \(n\)
\[
 \gcd\bigl(T_{2c+1,c,2}(n),\,c\bigr)=\gcd(n,c),
\]
and the same identity holds with \(T^{\,k}(n)\) in place of \(T(n)\).
\end{lemma}

\begin{proof}
If \(2\mid n\), then \(T(n)=n/2\) and \(n=2(n/2)\); since \(c\) is odd, the
factor \(2\) is coprime to \(c\), so \(\gcd(n/2,c)=\gcd(2(n/2),c)=\gcd(n,c)\).
If \(2\nmid n\), then \(T(n)=(2c+1)n+c=n+c(2n+1)\), and
\(\gcd(n+c(2n+1),c)=\gcd(n,c)\) because adding a multiple of \(c\) does not
change the gcd with \(c\). The iterate statement follows by induction on \(k\).

\end{proof}

\begin{theorem}[Odd offsets cycle inside their stratum]\label{Thm:StratumCycle}
Let \(c\) be odd and suppose the orbit of \(n\) is bounded. Then there are
\(y\) and \(p>0\) with
\[
 \gcd(y,c)=\gcd(n,c),\qquad T_{2c+1,c,2}^{\,p}(y)=y .
\]
\end{theorem}

\begin{proof}
By Lemma~\ref{Lem:Pigeonhole} there are \(k\) and \(p>0\) with
\(T^{\,k+p}(n)=T^{\,k}(n)\). Set \(y=T^{\,k}(n)\). The iterate identity of
Lemma~\ref{Lem:OddGcd} gives \(\gcd(y,c)=\gcd(n,c)\), and
\(T^{\,p}(y)=T^{\,k+p}(n)=T^{\,k}(n)=y\), using the addition rule to move the
first \(k\) instructions.
\end{proof}

\begin{corollary}[All strata carry cycles]\label{Cor:ManyCycles}
Let \(c\) be odd and positive, and suppose that no positive seed of the slice
\(T_{2c+1,c,2}\) diverges. Then for every divisor \(g\) of \(c\) there are
\(y\) and \(p>0\) with
\[
 \gcd(y,c)=g,\qquad T_{2c+1,c,2}^{\,p}(y)=y .
\]
\end{corollary}

\begin{proof}
A divisor of \(c\) is bounded by \(c\), so every seed is bounded under the
hypothesis, in particular the seed \(g\) itself, which is positive because
\(g\mid c\) and \(c>0\). Theorem~\ref{Thm:StratumCycle} applied to \(n=g\)
gives a periodic \(y\) with \(\gcd(y,c)=\gcd(g,c)=g\), the last identity
because \(g\) divides \(c\).
\end{proof}

\begin{corollary}[The seed of the offset never reaches \(1\)]\label{Cor:NoSeedCapture}
Let \(c\) be odd with \(c>1\). Then \(T_{2c+1,c,2}^{\,k}(c)\ne1\) for every
\(k\ge0\).
\end{corollary}

\begin{proof}
By Lemma~\ref{Lem:OddGcd} with \(n=c\),
\[
 \gcd\bigl(T_{2c+1,c,2}^{\,k}(c),\,c\bigr)=\gcd(c,c)=c .
\]
If \(T_{2c+1,c,2}^{\,k}(c)=1\) held for some \(k\), the left side would equal
\(\gcd(1,c)=1\), hence \(c=1\), contrary to the hypothesis.
\end{proof}

\section{The geometric-offset slice}

\begin{lemma}[The offset is integral]\label{Lem:GeoOffset}
For every \(h\), \(\;7\mid 8^{h}-1\), and the number
\(c_h=3(8^{h}-1)/7\) satisfies
\[
 7c_h=3\bigl(8^{h}-1\bigr),\qquad
 3q_h+c_h=3\cdot2^{3h}\quad\text{for } q_h=2c_h+1 .
\]
\end{lemma}

\begin{proof}
Since \(8\equiv1\pmod7\) we have \(8^{h}\equiv1\pmod7\) for every \(h\), so
\(7\mid8^{h}-1\); hence \(7\mid3(8^{h}-1)\) and \(c_h\) is an integer with
\(7c_h=3(8^{h}-1)\). Then
\(3q_h+c_h=3(2c_h+1)+c_h=7c_h+3=3(8^{h}-1)+3=3\cdot8^{h}\), and
\(8^{h}=2^{3h}\).
\end{proof}

\begin{theorem}[The cycle of the geometric slice]\label{Thm:GeoLoop}
On the slice \(d=2\), \(q=2c_h+1\), \(c=c_h\) of Lemma~\ref{Lem:GeoOffset},
\[
 T_{q_h,c_h,2}^{\,j+1}(3)=3\cdot2^{\,3h-j}\quad(0\le j\le3h),
\]
so \(T^{\,3h+1}(3)=3\), no positive time smaller than \(3h+1\) returns \(3\)
to itself, and the orbit meets \(1\) at no time.
\end{theorem}

\begin{proof}
The seed \(3\) is odd, so by Lemma~\ref{Lem:TwoOdd}
\(T(3)=3q_h+c_h=3\cdot2^{3h}\), which is the formula for \(j=0\). For the
step, let \(0\le j<3h\) and assume \(T^{\,j+1}(3)=3\cdot2^{\,3h-j}\). Since
\(2^{\,3h-j}\) is even, that value is divisible by \(2\), so one more
instruction divides it by \(2\):
\[
 T^{\,j+2}(3)=\frac{3\cdot2^{\,3h-j}}2=3\cdot2^{\,3h-j-1}
 =3\cdot2^{\,3h-(j+1)},
\]
which is the formula at \(j+1\). This proves the formula by induction.

At \(j=3h\) the formula gives \(3\cdot2^{0}=3\), so \(3\) returns. If
\(0<i<3h+1\), write \(i=j+1\) with \(0\le j<3h\); the formula gives
\(T^{\,i}(3)=3\cdot2^{\,e}\) with \(e=3h-j\ge1\), so
\(T^{\,i}(3)\ge6\ne3\). Finally, every value \(3\cdot2^{\,e}\) with \(e\ge0\)
is at least \(3\), so the orbit never meets \(1\): times larger than
\(3h+1\) reduce to times in \(\{0,\dots,3h\}\) by removing full periods of
length \(3h+1\), since \(T^{\,3h+1}(3)=3\).
\end{proof}

\section{The full-residue-access slice}

\begin{lemma}[Slice identities]\label{Lem:FullIdentity}
Let \(d\ge2\), \(k\ge2\) and set \(c=(d-1)(d^{k-1}-1)\),
\(q=(d-1)d^{k-1}+1\). Then
\[
 q=d+c,\qquad q(d-1)+c=(d-1)d^{k},\qquad d\nmid d-1 .
\]
\end{lemma}

\begin{proof}
The first identity is
\(q-c=(d-1)d^{k-1}+1-(d-1)d^{k-1}+(d-1)=d\). For the second, multiply out:
\[
 q(d-1)+c=\bigl((d-1)d^{k-1}+1\bigr)(d-1)+(d-1)\bigl(d^{k-1}-1\bigr)
 =(d-1)\Bigl(d^{k-1}(d-1)+1+d^{k-1}-1\Bigr)=(d-1)d^{k-1}d,
\]
and \(d^{k-1}d=d^{k}\) because \(k=(k-1)+1\). For the third, \(d\mid d-1\)
would give \(d\le d-1\) for \(d\ge2\), which is false.
\end{proof}

\begin{theorem}[The cycle of the full-residue-access slice]\label{Thm:FullLoop}
Let \(d\ge2\), \(k\ge2\) and take \(c\), \(q\) as in
Lemma~\ref{Lem:FullIdentity}. Then
\[
 T_{q,c,d}^{\,j+1}(d-1)=(d-1)d^{\,k-j}\quad(0\le j\le k),
\]
so \(T^{\,k+1}(d-1)=d-1\), no positive time smaller than \(k+1\) returns
\(d-1\) to itself, and for \(d\ge3\) the orbit meets \(1\) at no time.
\end{theorem}

\begin{proof}
By Lemma~\ref{Lem:FullIdentity} the state \(d-1\) is not divisible by \(d\), so
the first instruction is affine:
\(T(d-1)=q(d-1)+c=(d-1)d^{k}\). This is the formula for \(j=0\). For the step,
let \(0\le j<k\) and assume \(T^{\,j+1}(d-1)=(d-1)d^{\,k-j}\). Since
\(k-j\ge1\), that value is divisible by \(d\), so one more instruction divides
it by \(d\), giving \((d-1)d^{\,k-j-1}=(d-1)d^{\,k-(j+1)}\). This proves the
formula by induction.

At \(j=k\) the formula gives \((d-1)d^{0}=d-1\), so \(d-1\) returns. If
\(0<i<k+1\), write \(i=j+1\) with \(0\le j<k\); then
\(T^{\,i}(d-1)=(d-1)d^{\,e}\) with \(e=k-j\ge1\), and since \(d\ge2\),
\(d^{\,e}\ge2\) and \(d-1\ge1\), so \((d-1)d^{\,e}\ge2>d-1\): no earlier
positive time returns. Finally, for \(d\ge3\) every value
\((d-1)d^{\,e}\) with \(e\ge0\) is at least \(2\), so no time meets \(1\),
again after removing full periods of length \(k+1\).
\end{proof}

\begin{corollary}[The classical cycle]\label{Cor:Classical}
At \(d=k=2\) the slice is \(T_{3,1,2}(n)=n/2\) for even \(n\) and
\(T_{3,1,2}(n)=3n+1\) for odd \(n\), and
\[
 T_{3,1,2}(1)=4,\qquad T_{3,1,2}^{\,3}(1)=1 .
\]
\end{corollary}

\begin{proof}
At \(d=k=2\) we have \(c=(2-1)(2^{1}-1)=1\) and \(q=(2-1)2^{1}+1=3\), so the
slice is \(n/2\) on even states and \(3n+1\) on odd states by
Lemma~\ref{Lem:Branch}. Then \(1\) is odd, \(3\cdot1+1=4\);
\(4\) is even, \(4/2=2\); and \(2\) is even, \(2/2=1\).
\end{proof}

\begin{remark}[Full residue access does not single out the classical target]
\label{Rem:NotClassicalOnly}
The slice of Theorem~\ref{Thm:FullLoop} has \(q\equiv c\equiv1\pmod d\), so
its affine residues range over every class modulo \(d\). Nevertheless the
cycle it carries is \(d-1\mapsto(d-1)d^{k}\mapsto\dots\mapsto d(d-1)\mapsto d-1\),
which at \(d=k=2\) is \(1\mapsto4\mapsto2\mapsto1\) by
Corollary~\ref{Cor:Classical} but for \(d\ge3\) avoids \(1\) by
Theorem~\ref{Thm:FullLoop}. Full residue access therefore does not force the
classical target; the two slices have different targets.
\end{remark}

\section{The intersection slice}

\begin{lemma}[Predecessor]\label{Lem:IntersectionPred}
Let \(d\ge3\) and suppose \(T_{d+1,d-1,d}(x)=d^{m}\). Then \(x=d^{\,m+1}\).
\end{lemma}

\begin{proof}
If \(d\mid x\), the division branch gives \(x/d=d^{m}\), that is
\(x=d^{\,m+1}\).

Suppose instead that \(d\nmid x\). The affine branch gives
\((d+1)x+(d-1)=d^{m}\). Reduce this identity modulo \(d+1\): since
\(d\equiv-1\pmod{d+1}\), the term \((d+1)x\) vanishes and the right side is
\((-1)^{m}\), so
\[
 (-1)^{m}\equiv d-1\equiv-2\pmod{d+1} .
\]
If \(m\) is even this reads \(1\equiv-2\), that is \(3\equiv0\pmod{d+1}\),
so \(d+1\mid3\) and hence \(d+1\le3\), which contradicts \(d\ge3\). If
\(m\) is odd it reads \(-1\equiv-2\), that is \(1\equiv0\pmod{d+1}\), so
\(d+1\mid1\) and hence \(d+1\le1\), which contradicts \(d\ge3\) as well.
Both parities are impossible, so the affine case does not occur and the first
case applies.
\end{proof}

\begin{theorem}[Iterates of the intersection slice]\label{Thm:IntersectionIterate}
Let \(d\ge3\) and \(x\ge1\). If \(T_{d+1,d-1,d}^{\,k}(x)=d^{\,t}\), then
\[
 x=d^{\,k+t} .
\]
\end{theorem}

\begin{proof}
Induction on \(k\). For \(k=0\) the hypothesis is \(x=d^{\,t}\), which is the
claim. For the step, let \(y=T^{\,k}(x)\); then
\(T(y)=T^{\,k+1}(x)=d^{\,t}\), so Lemma~\ref{Lem:IntersectionPred} gives
\(y=d^{\,t+1}\). The induction hypothesis applied to \(y\) yields
\(x=d^{\,k+(t+1)}=d^{\,k+1+t}\).
\end{proof}

\begin{theorem}[The seeds that reach \(1\)]\label{Thm:IntersectionReach}
Let \(d\ge3\) and \(x\ge1\). Then some iterate of \(x\) equals \(1\) if and
only if \(x\) is a power of \(d\); for each \(k\),
\[
 T_{d+1,d-1,d}^{\,k}(x)=1 \iff x=d^{\,k} .
\]
\end{theorem}

\begin{proof}
If \(T^{\,k}(x)=1=d^{0}\), then Theorem~\ref{Thm:IntersectionIterate} with
\(t=0\) gives \(x=d^{\,k}\). Conversely, if \(x=d^{\,j}\), then
Lemma~\ref{Lem:Divisions} gives \(T^{\,j}(x)=1\), since \(d\ge2\).
\end{proof}

\begin{corollary}[No return of \(1\)]\label{Cor:NoReturn}
Let \(d\ge3\). Then \(T_{d+1,d-1,d}^{\,k}(1)\ne1\) for every \(k>0\).
\end{corollary}

\begin{proof}
If \(T^{\,k}(1)=1\) with \(k>0\), then
Theorem~\ref{Thm:IntersectionReach} gives \(1=d^{\,k}\), that is
\(d^{\,k}=d^{0}\); injectivity of the exponent for \(d\ge2\) forces \(k=0\),
contradicting \(k>0\).
\end{proof}

\section{Controls and excluded cases}

\begin{remark}[The divisor two is necessary]\label{Rem:TwoNecessary}
At \(d=2\) the intersection slice is \(q=3\), \(c=1\), the classical map. The
seed \(3\) reaches \(1\) through
\[
 3\mapsto10\mapsto5\mapsto16\mapsto8\mapsto4\mapsto2\mapsto1,
\]
and \(3\) is not a power of \(2\), so Theorem~\ref{Thm:IntersectionReach} fails
without \(d\ge3\). Likewise, at \(d=2\) the full-residue-access slice of
Theorem~\ref{Thm:FullLoop} has the seed \(1\) on the classical cycle
\(1\mapsto4\mapsto2\mapsto1\), so the hypothesis \(d\ge3\) in its last
assertion is exactly what excludes a meeting with \(1\).
\end{remark}

\begin{remark}[Zero and the odd-offset hypothesis]\label{Rem:Zero}
The state \(0\) is fixed by every raw map, by Lemma~\ref{Lem:Branch}, so
positivity is required wherever an orbit is claimed to grow. At \(d=2\),
\(q=1\), \(c=0\) the odd state \(1\) is fixed as well, since the affine branch
is \(1\mapsto1\cdot1+0\); the hypothesis \(1\le c\) in
Theorem~\ref{Thm:EvenEscape} excludes this degeneration. With odd \(c\) the
parity alternation of Theorem~\ref{Thm:EvenEscape} fails, because the affine
branch of an odd state is then even; the corresponding statement for odd
offsets is the stratum theorem, Theorem~\ref{Thm:StratumCycle}.
\end{remark}

\begin{remark}[Boundary parameters]\label{Rem:Boundary}
At \(h=0\) the geometric slice has vanishing offset and multiplier \(1\), and
the theorem still holds: the cycle of \(3\) has length \(1\). At \(k=1\) the
full-residue-access slice has vanishing offset as well, and the hypothesis
\(2\le k\) of Theorem~\ref{Thm:FullLoop} is required by the induction, which
divides \(k-j\) times.
\end{remark}

\section{Formal correspondence}

\begin{table}[htbp]
\centering
\begin{tabular}{ll}
\toprule
Statement & Lean declaration \\
\midrule
Lemma~\ref{Lem:Branch} & \texttt{RawStepOfDvd}, \texttt{RawStepOfNotDvd}, \texttt{RawStepZero} \\
Lemma~\ref{Lem:IterateSplit} & \texttt{RawStepIterateSucc}, \texttt{RawStepIterateSuccLast}, \texttt{IterateAdd} \\
Lemma~\ref{Lem:Divisions} & \texttt{RawStepDivMulIterate}, \texttt{RawStepPowIterate} \\
Lemma~\ref{Lem:Pigeonhole} & \texttt{RawStepBoundedEventuallyPeriodic} \\
Lemma~\ref{Lem:TwoOdd} & \texttt{RawStepTwoOfOdd}, \texttt{RawStepTwoEvenOffsetAffine}, \texttt{RawStepTwoEvenOffsetOddOfOdd}, \texttt{RawStepTwoEvenOffsetIncreaseOdd} \\
Theorem~\ref{Thm:EvenEscape} & \texttt{RawStepTwoEvenOffsetReachesOdd}, \texttt{RawStepTwoEvenOffsetUnbounded} \\
Lemma~\ref{Lem:OddGcd} & \texttt{RawStepTwoOddOffsetGcd}, \texttt{RawStepTwoOddOffsetGcdIterate} \\
Theorem~\ref{Thm:StratumCycle} & \texttt{RawStepTwoOddOffsetStratumCycle} \\
Corollary~\ref{Cor:ManyCycles} & \texttt{RawStepTwoOddOffsetManyCycles} \\
Corollary~\ref{Cor:NoSeedCapture} & \texttt{RawStepTwoOddOffsetSeedNoCapture} \\
Lemma~\ref{Lem:GeoOffset} & \texttt{SevenDvdEightPowSubOne}, \texttt{CycleC}, \texttt{CycleQ}, \texttt{CycleCMulSeven}, \texttt{CycleQIdentity} \\
Theorem~\ref{Thm:GeoLoop} & \texttt{CycleAffineStep}, \texttt{CycleIterateFormula}, \texttt{CycleReturns}, \texttt{CyclePrimitive}, \texttt{CycleAvoidsOne} \\
Lemma~\ref{Lem:FullIdentity} & \texttt{FullAccessQeq}, \texttt{FullAccessIdentity}, \texttt{FullAccessNotDvd}, \texttt{NatPowDvdSelfPow} \\
Theorem~\ref{Thm:FullLoop} & \texttt{FullAccessIterateFormula}, \texttt{FullAccessReturns}, \texttt{FullAccessPrimitive}, \texttt{FullAccessAvoidsOne} \\
Corollary~\ref{Cor:Classical} & \texttt{FullAccessClassicalInstance}, \texttt{FullAccessIterateFormula} \\
Lemma~\ref{Lem:IntersectionPred} & \texttt{RawStepIntersectionPredecessor} \\
Theorem~\ref{Thm:IntersectionIterate} & \texttt{RawStepIntersectionIteratePow} \\
Theorem~\ref{Thm:IntersectionReach} & \texttt{RawStepIntersectionReachesOne}, \texttt{RawStepIntersectionReachesOneIff}, \texttt{RawStepPowIterate} \\
Corollary~\ref{Cor:NoReturn} & \texttt{RawStepIntersectionNoReturnToOne} \\
Remark~\ref{Rem:NotClassicalOnly} & \texttt{FullAccessClassicalInstance}, \texttt{FullAccessAvoidsOne} \\
Remark~\ref{Rem:TwoNecessary} & \texttt{FullAccessClassicalInstance}, \texttt{RawStepIntersectionReachesOneIff} \\
Remark~\ref{Rem:Zero} & \texttt{RawStepZero}, \texttt{RawStepTwoOddOffsetGcd} \\
Remark~\ref{Rem:Boundary} & \texttt{CycleReturns}, \texttt{FullAccessIterateFormula} \\
\bottomrule
\end{tabular}
\caption{Lean declarations carrying each statement.}
\end{table}

\end{document}
